(1) is the nullity computation of Fact \claimref{one-relation.imported-pure-spinors-incidence-chart}. The script re-verifies it on the flat model with the inner product computed, not assumed: with $n_{0} = (e_{+} + e_{-})/2$, $n_{\infty} = e_{+} - e_{-}$ for $e_{\pm}^{2} = \pm 1$, it confirms $n_{0}^{2} = n_{\infty}^{2} = 0$, $\ip{n_{0}}{n_{\infty}} = 1$, then for six rational points with $Q(x)$ timelike, null and spacelike computes $Q(X(x)) = 0$ and $X(x)^{2} = 0$ in the algebra with $X(x) \neq 0$, and confirms that any other slot coefficient $c$ gives $Q(X) = Q(x)(1 + 2c) \neq 0$. A nonzero nilpotent is not invertible (\claimref{std-algebra.inverse-of-vector}: a vector is invertible exactly when $Q \neq 0$). (2) By Fact \claimref{one-relation.imported-pure-spinors-incidence-chart} a nonzero conformal spinor $u$ has $\dim W_{u} = 3$ and $W_{u}$ determines $[u]$; and $c(X)u = 0$ if and only if $X \in W_{u}$, by the definition of $W_{u}$. So the conformal spinors annihilated by $c(X(x))$ are exactly those whose isotropic system contains $X(x)$, and the set of such $[u]$ is the incidence line of the chart theorem, the line of $\mathbb{P}(\CS)$ that the simple bivector $X(x)$ represents. The kernel dimensions follow from nilpotency and the module structure: on the eight-dimensional Fock module a nonzero null vector has kernel of dimension exactly four, because the nonzero null vectors form a single $O(6,\C)$-orbit (Witt's extension theorem, Fact \claimref{one-relation.imported-witts-extension-theorem}), $c$ is equivariant under the Pin lift, so the rank of $c(X)$ is constant on the orbit, and at $X = w_{1}$ the operator is the wedge with one basis element of $W$, which annihilates exactly the four monomials containing it; and $c(X)$ swaps the chirality halves, so its kernel splits as two plus two. The script builds the Fock model explicitly (wedge and contraction on $\Lambda^{\bullet}W$, $\dim 8$, with the six ambient generators realised and all thirty-six anticommutators checked), and for each of the six points computes: $c(X)^{2} = 0$; kernel of dimension four, and two on each half; that two independent conformal spinors on the line and their sum each have a three-dimensional annihilator, pairwise distinct; that $X(x)$ lies in each of those annihilators; and, across the fifteen pairs of points, that the lines are pairwise distinct and meet exactly when the points are null-separated. The null-separation criterion in general: expanding with $n_{0}^{2} = n_{\infty}^{2} = 0$ and $\ip{n_{0}}{n_{\infty}} = 1$ gives $\ip{X(x)}{X(y)} = \ip{x}{y} - \tfrac12 Q(x) - \tfrac12 Q(y) = -\tfrac12 Q(x-y)$, which the script confirms on all fifteen pairs; and two incidence lines share a vacuum $u$ if and only if $X(x), X(y) \in W_{u}$ for some $u$, if and only if $X(x)$ and $X(y)$ span a totally isotropic plane (any two vectors of a maximal isotropic subspace are orthogonal; conversely an isotropic plane of $\C^{6}$ lies in exactly two maximal isotropic subspaces, which meet in it, so their intersection has odd codimension one in each and they lie in opposite families, exactly one of them the annihilator of a conformal spinor in $\CS$, which is also why two meeting incidence lines meet in one point), if and only if $\ip{X(x)}{X(y)} = 0$, since both are null. (3) and (4) are readings of (1)-(2) together with the imported theorems named (for (3) the CAR of \claimref{one-relation.two-presentations}; for (4) \claimref{sheaf.one-sided-completion} and \claimref{sheaf.localised-transport}, in the Schwarzschild form of \claimref{sheaf.one-sided-completion-singularity}), and are stated at that level: the content is the placement of the two theories on one object, not a new dynamical claim.
